% Generated from validated Article Draft Result. Do not hand-edit; revise the structured draft instead.

GPT-3後の規模拡大だけを見ていると、2021年に生じたもう一つの変化を見落とす。モデル本体を毎回全面的に作り直すのではなく、既存のpretrained modelを用途へ合わせる方法そのものが研究対象として分化したことである。本稿ではPrefix-Tuning、Prompt Tuning、LoRAを、同じ結論へ収束した三手法ではなく、再利用可能な基盤と個別タスクの間をどう接続するかという共通問題に対する別々の設計として読む。\autocite{src-ab9d63ddee6e,src-c1ad400391c9,src-dba095cd4e62}


Prefix-Tuningが示す設計上の焦点は、full fine-tuningと同じようにモデル全体を更新することではなく、prefix parametersを介して既存モデルを適応させる点にある。ここで重要なのは『小さいから同じ』ではないことで、更新対象をモデル本体から切り離すという境界が、用途別適応を一つの独立問題として見せた。性能評価や優位性については原著者の評価として扱い、本稿は手法の位置付けまでを年次記録として採用する。\autocite{src-dba095cd4e62}


Prompt Tuningは、soft promptを適応対象とする別の経路を示した。通常の自然言語promptを実行時に工夫することと、学習可能なsoft promptを通じてモデルを適応させることは分けて考える必要がある。2021年には、promptが単なる入力文の書き方ではなく、基盤モデルを再利用するための学習可能なinterfaceにもなり得るという設計空間が明瞭になった。\autocite{src-c1ad400391c9}


LoRAは、low-rank updatesによって更新を構成し、trainable parametersを抑えるという方向を代表する。Prefixやsoft promptとは更新対象・機構が異なるため、三者を一つの技術名でまとめると2021年の設計上の分化が消えてしまう。共通しているのは、巨大モデルを用途の数だけ複製・再学習する発想から、固定された基盤と小さな適応部分を組み合わせる発想へ重心を移した点である。\autocite{src-ab9d63ddee6e}


年次の視点では、三手法の勝敗よりも、pretrainingの後に独立したadaptation stageを置けることの方が重要である。Prefix parameters、soft prompts、low-rank updatesという異なる選択肢が並んだことで、『一つの巨大モデルを作る』工程と『目的に合わせて振る舞いを変える』工程を分離して設計できる余地が広がった。この分離は、同じ2021年に進んだinstruction tuningやhuman feedbackとも接続し、翌年以降のpost-trainingを理解するための前提になる。\autocite{src-ab9d63ddee6e,src-c1ad400391c9,src-dba095cd4e62}


\begin{claimboundary}
ここで確認しているのは、各手法が2021年の一次資料に記録され、異なるparameter-efficient adaptationの設計を提示したことまでである。原著者が報告した性能・効率上の評価を独立再現済みの事実へ置き換えない。\autocite{src-ab9d63ddee6e,src-c1ad400391c9,src-dba095cd4e62}
\end{claimboundary}
